\documentclass[11pt,letterpaper]{article}
\usepackage[T1]{fontenc}
\usepackage{lmodern}
\usepackage[margin=1.05in]{geometry}
\usepackage{amsmath,amssymb,amsthm,booktabs,longtable,microtype,xurl}
\usepackage[hidelinks,pdftitle={Conditional ZFC + I3 independence for a 334-state Turing machine},
  pdfauthor={James Chang}]{hyperref}

\newtheorem{theorem}{Theorem}[section]
\newtheorem{lemma}[theorem]{Lemma}
\newcommand{\ZFC}{\mathrm{ZFC}}
\newcommand{\Ithree}{\mathrm{I}_3}
\newcommand{\Con}{\operatorname{Con}}
\newcommand{\Halt}{\operatorname{Halt}}
\newcommand{\code}[1]{\texttt{\nolinkurl{#1}}}
\newcommand{\nand}{\mathbin{\uparrow}}
\newcommand{\conv}[1]{#1^{\smile}}
\newcommand{\ch}{\operatorname{ch}}
\setlength{\emergencystretch}{2em}

\title{Conditional \texorpdfstring{$\ZFC+\Ithree$}{ZFC + I3} independence\\
for a 334-state Turing machine}
\author{James Chang}
\date{}

\begin{document}
\maketitle

\begin{abstract}
We give a two-symbol Turing machine with 334 working states whose
blank-tape nonhalting is equivalent to the syntactic consistency of
$\ZFC+\Ithree$. The machine searches for an equality in a calculus of
ten guarded equations. Consistency suffices for the unprovability of
nonhalting, and arithmetic $\Sigma_1$-soundness suffices for the
unprovability of halting. The proof distinguishes standard finite
computations from internally finite computations in models of the
theory. Lean checks the literal transition table, the guarded search,
and the consistency implication needed for the second incompleteness
argument.
\end{abstract}

\section{Statement}

Write $T=\ZFC+\Ithree$. Here $\Ithree$ asserts that some $V_\lambda$,
for a nonzero limit ordinal $\lambda$, admits a nonidentity elementary
embedding into itself in the language of membership
\cite{corazza}. The first-order axiom expresses elementarity using
satisfaction for the set structure $V_\lambda$. In a model of $T$ this
satisfaction condition ranges over the formula codes of that model.

Let $M$ be the machine in Appendix~\ref{app:table}. It starts in
state~$0$ on a doubly infinite tape of zeros. States $0,\ldots,333$ are
working states, and $334$ is the halt label. The unique halting
transition reads $1$ at state~$189$, writes $1$, and moves right.
Write $\Halt(M)$ for actual halting and $H_M$ for the closed
set-theoretic sentence asserting that a finite halting history exists.

\begin{theorem}\label{thm:main}
In every model of $T$, $H_M$ holds exactly when there is a blank-tape
halting computation in that model. Moreover,
\begin{enumerate}
\item $\neg\Halt(M)\iff\Con(T)$.
\item If $T$ is syntactically consistent, then $T\nvdash\neg H_M$.
\item If $T$ is arithmetically $\Sigma_1$-sound, then $T\nvdash H_M$
and $T\nvdash\neg H_M$.
\end{enumerate}
\end{theorem}

Arithmetic $\Sigma_1$-soundness means that a closed arithmetic
$\Sigma_1$ sentence is true in $\mathbb N$ whenever $T$ proves a
set-theoretic sentence representing it. Representation here is
agreement in every model of $T$ with the arithmetic sentence
interpreted on the model's natural numbers.

\section{The guarded equations}

A symbol in the machine's term language is a pair $(k,t)$, where $k$
is the arity and $t$ is the tag. A term is read in prefix order. Use
$J,D,E,K,I$ for the constants with tags $0,1,2,3,5$, respectively.
The unary operations $\ch$, $\neg$ and converse have tags $1,3,4$.
NAND $\nand$ and product $\cdot$ have tags $3,4$ at arity two. Set
\[
x\wedge y=\neg(x\nand y),\qquad
U=(\neg E)\nand E,\qquad
P=\neg(\conv E\cdot\neg E),\qquad S=\conv{\ch(P)}\cdot\ch(P).
\]
The variables $a,b,c$ range over all closed terms. The variable $d$
ranges over closed terms in which every tag is at least three.
Write $F=\ch(a)\wedge a$ and $X=\conv E\cdot F$. The six factors of
the guarded equation are
\begin{align*}
B_1 &= \bigl(\neg((\neg X)\cdot E) \wedge
                 \neg(\conv E\cdot(F\cdot\neg E))\bigr)\cdot U,\\
B_2 &= (\conv a\cdot(\neg(a\cdot b)\wedge c))\nand b,\\
B_3 &= (d\cdot J)\nand(J\cdot\neg d),\\
B_4 &= a\nand\neg(F\cdot U),\\
B_5 &= (\conv{\ch(a)}\cdot\ch(a))\nand\neg I,\\
B_6 &= (\conv P)\nand(P\wedge\neg I).
\end{align*}
Using these abbreviations, the ten equations in compiled order are
\begin{align*}
(\conv a)\nand(\conv b)&=\conv{(a\nand b)},\\
a\cdot I&=a,\\
B_1\wedge(B_2\wedge(B_3\wedge(B_4\wedge(B_5\wedge B_6))))&=U,\\
P\cdot(S\cdot(E\cdot S))&=E\cdot S,\\
a&=\neg\neg((\neg a)\nand b)\wedge a,\\
a\cdot(b\cdot c)&=(a\cdot b)\cdot c,\\
I\wedge(J\cdot U)&=D,\\
b\wedge(c\wedge a)&=\neg\neg((\neg a)\nand b)\wedge(c\wedge b),\\
I\wedge(K\cdot(S\cdot U))&=D,\\
K&=E\cdot D.
\end{align*}
The machine allows either orientation of an instantiated equation in
a common prefix and suffix context. A capture of $d$ checks the guard
on the entire substituted term.

\begin{lemma}\label{lem:consistent}
The guarded calculus fails to derive $D=J$ exactly when $T$ is
consistent.
\end{lemma}

\begin{proof}
In a model of $T$ the algebra of definable relations interprets $D$ as
the diagonal of $V_\lambda$, $E$ as membership and $J$ as the graph of
the embedding. The Meredith equations give the Boolean reduct
\cite{robbins}. The remaining identities express composition,
choice, domain and stage laws. The seed $K=E\cdot D$ is preserved by
the embedding. The closure proof for guarded terms propagates
preservation through the allowed operations, including a complement
argument on the domain of the embedding. Hence the interpretation
satisfies the guarded equation. Nonidentity of the embedding gives
$D\ne J$, so soundness of the calculus yields underivability.

Conversely, form the ground-term algebra modulo guarded derivability.
Assuming $D\ne J$, recover the relation-algebra laws from the ten
equations. A guarded class is one with a representative satisfying
the tag bound. The seed is guarded, and the class of
$(x\cdot\conv y)\wedge(y\nand z)$ is guarded when $x,y,z$ are
guarded. To interpret the older local-relation symbol, use the product
with a two-element Boolean selector. A reserved selector value names
the seed, and an available guarded representative acts as the
local relation of its own class. This supplies the local closure and
preservation equations of the older calculus. The representation
argument for that calculus then produces a model of $T$.
\end{proof}

\section{From the literal tape to proof search}

The simulation first checks the startup trace and the layout of the
program and unary registers. The decoder is verified against a
344-state reference using a finite state correspondence and bounded
tape windows. The register proof handles growth, return traversal, and
the actual halt transition in the 334-state table. The structured
program read by the decoder normalizes to the register program for
the equations above.

A candidate input contains a proposed term, a context position and
orientation, four substitution words, and a sequence of selectors for
the equation trie. The register program compares and removes a common
prefix, interprets the requested equation, scans each occurrence of a
capture, and compares the remaining suffixes. A successful trial
replaces the incumbent term. The distinct codes $4$ and $1$ represent
$D$ and $J$, and the acceptance gate halts exactly when a checked
proposal has code~$1$.

For completeness, a single guarded rewrite determines both the
ordinary trie selectors and any branch selected by a dictionary
suffix. A finite proof is packed backwards into one input. Between
trials the scheduler subtracts one from the pending word, so each
continuing proof suffix includes this debit. All arbitrary trials
terminate or halt. Below a fixed target input $N$, the expression
$(N-i)(N+1)+p$ bounds the wait from index~$i$ and pending word~$p$.
It decreases when an attempt spends pending input or advances the
outer index. Every chosen finite proof therefore leads to a halting
run. The soundness of each accepted rewrite gives the reverse
implication. Together with Lemma~\ref{lem:consistent}, this proves the
external equivalence in Theorem~\ref{thm:main}.

\section{Internal calibration}

An implication about standard finite derivations does not suffice
to apply the second incompleteness theorem. The calibration uses a
recursive arithmetic theory $A$ extending $\mathrm{PA}$, whose
arithmetic consequences include the translated consequences of $T$
needed in the argument. Consistency of $T$ implies consistency of~$A$.
We show that~$T$ proves $\neg H_M\to\Con(A)$.

Work in the arithmetic of a model of $T$. If the older coded
ground-term quotient separates $D$ and $J$, an arithmetic-coded graph
and its truth table exclude an internal refutation of~$A$. Suppose
instead that the quotient collapses $D$ and $J$. In the arithmetic
model, form the native quotient by guarded rewrite traces of
internally finite length. The same guarded-class construction as in
Lemma~\ref{lem:consistent} validates each older equation in this
quotient. To transport a coded old derivation, evaluate its prefix
words by a right-to-left stack fold. Its graph is arithmetically
definable, including the test for a guarded representative. The
evaluation preserves equality under concatenation and the two-sided
rewrite contexts. Arithmetic induction through the coded old trace
then yields a native guarded derivation of $D=J$.

The native derivation is replayed by the register program with all
word lengths, queue debits and loop counts interpreted in the same
arithmetic model. The checked entry frame and decoder simulation turn
that run into a physical halting history of the exact table. Thus
$\neg H_M$ excludes the collapse and the refutation, which proves
the desired consistency implication by first-order completeness.

If $T$ proved $\neg H_M$, it would prove the interpretation of
$\Con(A)$, so $A$ would prove its own consistency. Under consistency
of $T$, this contradicts the second incompleteness theorem
\cite{goedel}. For the remaining direction, arithmetic halting for
the exact table is $\Sigma_1$. The representation of $H_M$ and the
external equivalence exclude a proof of halting under arithmetic
$\Sigma_1$-soundness.

\section{Formal verification}

The Lean declaration \code{PalomarI3334.exactClaim} gives the literal
conditional independence statement. The exact external equivalence
is recorded separately as
\code{PalomarI3334.literal_nonhalting_iff_consistency}.
The proof uses the arithmetic and incompleteness library Foundation
\cite{foundation}, the binary tape semantics and proof transport adapted from
the 380-state formalization, and a standalone Challenge
compared against the Solution. The accompanying reproduction notes
give the build, axiom audit and Comparator commands.

\paragraph{Acknowledgment.}
Codex assisted with the Lean proof, verification and paper preparation.

\begin{thebibliography}{9}
\bibitem{corazza}
P. Corazza, The gap between $\mathrm I_3$ and the wholeness axiom,
\emph{Fundamenta Mathematicae} \textbf{179} (2003), 43--60.
\href{https://doi.org/10.4064/fm179-1-4}{doi:10.4064/fm179-1-4}.

\bibitem{robbins}
W. Truszkowska and A. Grabowski, On the two short axiomatizations of
ortholattices, \emph{Journal of Formalized Mathematics} \textbf{15}
(2003), Mizar article ROBBINS2.
\url{https://mizar.uwb.edu.pl/JFM/Vol15/robbins2.html}.

\bibitem{goedel}
K. G\"odel, \"Uber formal unentscheidbare S\"atze der Principia
Mathematica und verwandter Systeme I, \emph{Monatshefte f\"ur
Mathematik und Physik} \textbf{38} (1931), 173--198.
\href{https://doi.org/10.1007/BF01700692}{doi:10.1007/BF01700692}.

\bibitem{foundation}
S. Saitou and M. Noguchi, \emph{Formalized Formal Logic}, Lean library.
\url{https://github.com/FormalizedFormalLogic/Foundation}.
\end{thebibliography}

\appendix
\section{Exact transition table}\label{app:table}

In each group the columns give the working state $q$ and the two
actions for input bits $0$ and $1$. Each action encodes the written
bit, a direction, and the successor state. The SHA-256 digest of the
published table is
\begin{center}\small
\code{50b86af5fde24402b7e2326171deff3b781c66b799f0a9854090f23edf68399d}.
\end{center}
\begingroup
\small
\setlength{\tabcolsep}{3pt}
\begin{longtable}{rll@{\qquad}rll@{\qquad}rll@{\qquad}rll}
\toprule
$q$ & $0$ & $1$ & $q$ & $0$ & $1$ & $q$ & $0$ & $1$ & $q$ & $0$ & $1$ \\
\midrule
\endhead
0 & \texttt{1R001} & \texttt{0L002} & 84 & \texttt{1L046} & \texttt{1R066} & 168 & \texttt{0R090} & \texttt{1R183} & 252 & \texttt{0R264} & \texttt{1R265} \\
1 & \texttt{0R002} & \texttt{0R001} & 85 & \texttt{0R031} & \texttt{1R090} & 169 & \texttt{0R152} & \texttt{1R184} & 253 & \texttt{0R205} & \texttt{1R215} \\
2 & \texttt{1L003} & \texttt{0R004} & 86 & \texttt{0R031} & \texttt{1R073} & 170 & \texttt{0R092} & \texttt{1R185} & 254 & \texttt{0R176} & \texttt{1R266} \\
3 & \texttt{0L004} & \texttt{0L005} & 87 & \texttt{0R041} & \texttt{1L099} & 171 & \texttt{0R186} & \texttt{0L008} & 255 & \texttt{0R176} & \texttt{1R139} \\
4 & \texttt{1L000} & \texttt{0L001} & 88 & \texttt{0L100} & \texttt{1R101} & 172 & \texttt{0R182} & \texttt{1R187} & 256 & \texttt{0R267} & \texttt{1R267} \\
5 & \texttt{0L005} & \texttt{0L006} & 89 & \texttt{0R031} & \texttt{1R067} & 173 & \texttt{0R188} & \texttt{1R112} & 257 & \texttt{0L268} & \texttt{1R243} \\
6 & \texttt{0L006} & \texttt{0L007} & 90 & \texttt{0R031} & \texttt{1R055} & 174 & \texttt{0R189} & \texttt{1R190} & 258 & \texttt{0R260} & \texttt{1R269} \\
7 & \texttt{0L007} & \texttt{0L008} & 91 & \texttt{0R075} & \texttt{1R102} & 175 & \texttt{0R191} & \texttt{1R192} & 259 & \texttt{0R092} & \texttt{1R270} \\
8 & \texttt{0L008} & \texttt{0L009} & 92 & \texttt{0R103} & \texttt{1R104} & 176 & \texttt{0R149} & \texttt{1R151} & 260 & \texttt{0R152} & \texttt{1R271} \\
9 & \texttt{0L010} & \texttt{1L009} & 93 & \texttt{0R043} & \texttt{1R105} & 177 & \texttt{0R075} & \texttt{1R193} & 261 & \texttt{0R140} & \texttt{1R272} \\
10 & \texttt{0L011} & \texttt{1L009} & 94 & \texttt{0R106} & \texttt{1R107} & 178 & \texttt{0R140} & \texttt{1R194} & 262 & \texttt{0R249} & \texttt{1R169} \\
11 & \texttt{0L012} & \texttt{1L009} & 95 & \texttt{1L108} & \texttt{0L046} & 179 & \texttt{0R195} & \texttt{1R136} & 263 & \texttt{0R273} & \texttt{1R140} \\
12 & \texttt{0L013} & \texttt{1L009} & 96 & \texttt{0R109} & \texttt{0R081} & 180 & \texttt{0R056} & \texttt{0L009} & 264 & \texttt{0R075} & \texttt{1R274} \\
13 & \texttt{0L014} & \texttt{1L009} & 97 & \texttt{0L110} & \texttt{1L082} & 181 & \texttt{0R196} & \texttt{1R197} & 265 & \texttt{0R220} & \texttt{1R275} \\
14 & \texttt{0R015} & \texttt{1L009} & 98 & \texttt{0L031} & \texttt{1R109} & 182 & \texttt{0R123} & \texttt{0L009} & 266 & \texttt{0R229} & \texttt{1R276} \\
15 & \texttt{0R015} & \texttt{1R016} & 99 & \texttt{0R111} & \texttt{1R112} & 183 & \texttt{0R075} & \texttt{1R198} & 267 & \texttt{0R273} & \texttt{1R217} \\
16 & \texttt{0R017} & \texttt{1R018} & 100 & \texttt{0R113} & \texttt{1R114} & 184 & \texttt{0R199} & \texttt{1R153} & 268 & \texttt{0R277} & \texttt{1R278} \\
17 & \texttt{0R019} & \texttt{1R020} & 101 & \texttt{0R115} & \texttt{0L008} & 185 & \texttt{0R139} & \texttt{1R200} & 269 & \texttt{0R092} & \texttt{1R279} \\
18 & \texttt{0R021} & \texttt{1R022} & 102 & \texttt{0R116} & \texttt{0L007} & 186 & \texttt{0R111} & \texttt{1R187} & 270 & \texttt{0R257} & \texttt{1R280} \\
19 & \texttt{0R023} & \texttt{0L024} & 103 & \texttt{0R067} & \texttt{1R117} & 187 & \texttt{0R089} & \texttt{0L009} & 271 & \texttt{0R199} & \texttt{1R184} \\
20 & \texttt{0R025} & \texttt{1R026} & 104 & \texttt{0R118} & \texttt{1R114} & 188 & \texttt{0R054} & \texttt{1R201} & 272 & \texttt{0R092} & \texttt{1R281} \\
21 & \texttt{0R027} & \texttt{1R028} & 105 & \texttt{0R119} & \texttt{1R120} & 189 & \texttt{0R202} & \texttt{1R334} & 273 & \texttt{0R150} & \texttt{1R164} \\
22 & \texttt{0R029} & \texttt{1R030} & 106 & \texttt{0R075} & \texttt{0L009} & 190 & \texttt{0R068} & \texttt{1R203} & 274 & \texttt{0R282} & \texttt{0L024} \\
23 & \texttt{0R031} & \texttt{1R032} & 107 & \texttt{0R121} & \texttt{1R122} & 191 & \texttt{0R075} & \texttt{1R204} & 275 & \texttt{0R162} & \texttt{1R283} \\
24 & \texttt{0L033} & \texttt{0L024} & 108 & \texttt{0R095} & \texttt{1L108} & 192 & \texttt{0R205} & \texttt{1R206} & 276 & \texttt{0R239} & \texttt{1R284} \\
25 & \texttt{0R034} & \texttt{1R035} & 109 & \texttt{0L047} & \texttt{1R109} & 193 & \texttt{0R207} & \texttt{0L024} & 277 & \texttt{0R176} & \texttt{1R239} \\
26 & \texttt{0R036} & \texttt{1R037} & 110 & \texttt{0L110} & \texttt{0L024} & 194 & \texttt{0R208} & \texttt{1R209} & 278 & \texttt{0R230} & \texttt{1R285} \\
27 & \texttt{0R023} & \texttt{1R038} & 111 & \texttt{0R123} & \texttt{0L024} & 195 & \texttt{0R181} & \texttt{1R210} & 279 & \texttt{0R176} & \texttt{1R286} \\
28 & \texttt{0R039} & \texttt{1R040} & 112 & \texttt{0R039} & \texttt{1R124} & 196 & \texttt{0R056} & \texttt{1R125} & 280 & \texttt{0R164} & \texttt{1L243} \\
29 & \texttt{0R041} & \texttt{1R042} & 113 & \texttt{0R071} & \texttt{1R125} & 197 & \texttt{0L145} & \texttt{1R211} & 281 & \texttt{0R253} & \texttt{1R287} \\
30 & \texttt{0R043} & \texttt{1R044} & 114 & \texttt{0R039} & \texttt{1R126} & 198 & \texttt{0R179} & \texttt{0L008} & 282 & \texttt{0R169} & \texttt{1R238} \\
31 & \texttt{1R045} & \texttt{0L046} & 115 & \texttt{0R031} & \texttt{1R127} & 199 & \texttt{0R196} & \texttt{1R212} & 283 & \texttt{0R288} & \texttt{1R289} \\
32 & \texttt{0R031} & \texttt{1R047} & 116 & \texttt{0R128} & \texttt{1R128} & 200 & \texttt{0R138} & \texttt{1R213} & 284 & \texttt{0R243} & \texttt{1R290} \\
33 & \texttt{1L009} & \texttt{0L033} & 117 & \texttt{0R069} & \texttt{0L008} & 201 & \texttt{0R211} & \texttt{0L008} & 285 & \texttt{0R291} & \texttt{1R292} \\
34 & \texttt{0R032} & \texttt{1R048} & 118 & \texttt{0L100} & \texttt{1R129} & 202 & \texttt{0R089} & \texttt{0L024} & 286 & \texttt{0R176} & \texttt{1R293} \\
35 & \texttt{0R039} & \texttt{1R049} & 119 & \texttt{0R075} & \texttt{1R130} & 203 & \texttt{0R089} & \texttt{1R129} & 287 & \texttt{0R253} & \texttt{1R213} \\
36 & \texttt{0R031} & \texttt{1R050} & 120 & \texttt{0R131} & \texttt{1R132} & 204 & \texttt{0R214} & \texttt{0L024} & 288 & \texttt{0R075} & \texttt{1R294} \\
37 & \texttt{0R051} & \texttt{1R052} & 121 & \texttt{0R067} & \texttt{1R133} & 205 & \texttt{0R215} & \texttt{1R138} & 289 & \texttt{0R280} & \texttt{1R295} \\
38 & \texttt{0R053} & \texttt{0L008} & 122 & \texttt{0R088} & \texttt{1R134} & 206 & \texttt{0R043} & \texttt{1R216} & 290 & \texttt{0R229} & \texttt{1R268} \\
39 & \texttt{0R031} & \texttt{1R054} & 123 & \texttt{0R031} & \texttt{1R041} & 207 & \texttt{0R217} & \texttt{1R218} & 291 & \texttt{0R296} & \texttt{1R297} \\
40 & \texttt{0R055} & \texttt{0L008} & 124 & \texttt{0R086} & \texttt{0L008} & 208 & \texttt{0R056} & \texttt{1R219} & 292 & \texttt{0R140} & \texttt{1R298} \\
41 & \texttt{0R031} & \texttt{1R056} & 125 & \texttt{0R135} & \texttt{0L008} & 209 & \texttt{0R220} & \texttt{1R221} & 293 & \texttt{0R213} & \texttt{1R260} \\
42 & \texttt{0R057} & \texttt{0L008} & 126 & \texttt{0R073} & \texttt{0L008} & 210 & \texttt{0R222} & \texttt{1R223} & 294 & \texttt{0R299} & \texttt{0L024} \\
43 & \texttt{0R021} & \texttt{1R058} & 127 & \texttt{0R031} & \texttt{1R069} & 211 & \texttt{0R127} & \texttt{1L048} & 295 & \texttt{0R296} & \texttt{1R300} \\
44 & \texttt{0R059} & \texttt{1R060} & 128 & \texttt{0R100} & \texttt{1R136} & 212 & \texttt{0R197} & \texttt{1R157} & 296 & \texttt{0R138} & \texttt{1R257} \\
45 & \texttt{0R045} & \texttt{1R061} & 129 & \texttt{0R085} & \texttt{0L008} & 213 & \texttt{0R140} & \texttt{1R169} & 297 & \texttt{0R225} & \texttt{1R255} \\
46 & \texttt{0L062} & \texttt{1L046} & 130 & \texttt{0R137} & \texttt{0L024} & 214 & \texttt{0R224} & \texttt{1R225} & 298 & \texttt{0R253} & \texttt{1R301} \\
47 & \texttt{1R063} & \texttt{1L064} & 131 & \texttt{0R138} & \texttt{1R139} & 215 & \texttt{0R142} & \texttt{1R226} & 299 & \texttt{0R226} & \texttt{1R302} \\
48 & \texttt{0R065} & \texttt{0R008} & 132 & \texttt{0R140} & \texttt{1R141} & 216 & \texttt{0R227} & \texttt{1R228} & 300 & \texttt{0R092} & \texttt{1R280} \\
49 & \texttt{0R036} & \texttt{0L008} & 133 & \texttt{0R142} & \texttt{0L008} & 217 & \texttt{0R229} & \texttt{1R230} & 301 & \texttt{0R286} & \texttt{1R303} \\
50 & \texttt{0R050} & \texttt{1R066} & 134 & \texttt{0R143} & \texttt{1R144} & 218 & \texttt{0R231} & \texttt{1R232} & 302 & \texttt{0R231} & \texttt{1R295} \\
51 & \texttt{0R067} & \texttt{0L009} & 135 & \texttt{0R069} & \texttt{1R145} & 219 & \texttt{0R233} & \texttt{0L024} & 303 & \texttt{0R169} & \texttt{1R304} \\
52 & \texttt{0R068} & \texttt{1R068} & 136 & \texttt{0R146} & \texttt{1R143} & 220 & \texttt{0R164} & \texttt{1R234} & 304 & \texttt{0R225} & \texttt{1R305} \\
53 & \texttt{0R069} & \texttt{1R070} & 137 & \texttt{0R147} & \texttt{1R148} & 221 & \texttt{0R092} & \texttt{1R140} & 305 & \texttt{0R176} & \texttt{1R306} \\
54 & \texttt{0R031} & \texttt{1R071} & 138 & \texttt{0R149} & \texttt{1R149} & 222 & \texttt{0R123} & \texttt{1R235} & 306 & \texttt{0R307} & \texttt{1R308} \\
55 & \texttt{0R031} & \texttt{1R036} & 139 & \texttt{0R150} & \texttt{1R151} & 223 & \texttt{0R159} & \texttt{1R236} & 307 & \texttt{0R152} & \texttt{1R309} \\
56 & \texttt{0R031} & \texttt{1R039} & 140 & \texttt{0R152} & \texttt{1R153} & 224 & \texttt{0R075} & \texttt{1R237} & 308 & \texttt{0R280} & \texttt{1R310} \\
57 & \texttt{0R072} & \texttt{1R073} & 141 & \texttt{0R139} & \texttt{1R154} & 225 & \texttt{0L238} & \texttt{1L239} & 309 & \texttt{0R311} & \texttt{1R312} \\
58 & \texttt{0R041} & \texttt{1R074} & 142 & \texttt{0R031} & \texttt{1R115} & 226 & \texttt{0R164} & \texttt{1R149} & 310 & \texttt{0R176} & \texttt{1R313} \\
59 & \texttt{0R075} & \texttt{1R076} & 143 & \texttt{0R155} & \texttt{1R156} & 227 & \texttt{0R075} & \texttt{1R240} & 311 & \texttt{0R314} & \texttt{1R199} \\
60 & \texttt{0R077} & \texttt{1R078} & 144 & \texttt{0R157} & \texttt{1R158} & 228 & \texttt{0R176} & \texttt{1R241} & 312 & \texttt{0R090} & \texttt{1R315} \\
61 & \texttt{0R079} & \texttt{1R061} & 145 & \texttt{0R111} & \texttt{1R159} & 229 & \texttt{0R215} & \texttt{1R242} & 313 & \texttt{0R280} & \texttt{1R316} \\
62 & \texttt{0L080} & \texttt{1L046} & 146 & \texttt{0R160} & \texttt{1R118} & 230 & \texttt{0R215} & \texttt{1R243} & 314 & \texttt{0R199} & \texttt{1R199} \\
63 & \texttt{0R063} & \texttt{0R081} & 147 & \texttt{0R075} & \texttt{1R161} & 231 & \texttt{0R176} & \texttt{1R244} & 315 & \texttt{0R075} & \texttt{1R317} \\
64 & \texttt{0L082} & \texttt{1R083} & 148 & \texttt{0R162} & \texttt{1R163} & 232 & \texttt{0R245} & \texttt{1R246} & 316 & \texttt{0R215} & \texttt{1R318} \\
65 & \texttt{0R055} & \texttt{1R069} & 149 & \texttt{0R142} & \texttt{1R150} & 233 & \texttt{0R176} & \texttt{1R247} & 317 & \texttt{0R151} & \texttt{0L008} \\
66 & \texttt{0R084} & \texttt{1R066} & 150 & \texttt{0R142} & \texttt{1R164} & 234 & \texttt{0R142} & \texttt{1R151} & 318 & \texttt{0R307} & \texttt{1R319} \\
67 & \texttt{0R031} & \texttt{1R075} & 151 & \texttt{0L165} & \texttt{1R149} & 235 & \texttt{0R074} & \texttt{0L008} & 319 & \texttt{0R225} & \texttt{1R320} \\
68 & \texttt{0R085} & \texttt{1R085} & 152 & \texttt{0R106} & \texttt{1R166} & 236 & \texttt{0R164} & \texttt{1R181} & 320 & \texttt{0R140} & \texttt{1R321} \\
69 & \texttt{0R031} & \texttt{1R086} & 153 & \texttt{0R167} & \texttt{1R168} & 237 & \texttt{0R248} & \texttt{0L024} & 321 & \texttt{0R215} & \texttt{1R322} \\
70 & \texttt{0R087} & \texttt{1R088} & 154 & \texttt{0R169} & \texttt{1R170} & 238 & \texttt{0R249} & \texttt{1R250} & 322 & \texttt{0R277} & \texttt{1R323} \\
71 & \texttt{0R031} & \texttt{1R089} & 155 & \texttt{0R041} & \texttt{1R171} & 239 & \texttt{0R205} & \texttt{1R251} & 323 & \texttt{0R291} & \texttt{1R324} \\
72 & \texttt{0R031} & \texttt{1R085} & 156 & \texttt{0R123} & \texttt{1R172} & 240 & \texttt{0R252} & \texttt{0L024} & 324 & \texttt{0R263} & \texttt{1R325} \\
73 & \texttt{0R031} & \texttt{1R072} & 157 & \texttt{0L173} & \texttt{0L009} & 241 & \texttt{0R253} & \texttt{1R254} & 325 & \texttt{0R263} & \texttt{1R326} \\
74 & \texttt{0R090} & \texttt{1L048} & 158 & \texttt{0R089} & \texttt{1R174} & 242 & \texttt{0R255} & \texttt{1R243} & 326 & \texttt{0R327} & \texttt{1R328} \\
75 & \texttt{0R031} & \texttt{1R023} & 159 & \texttt{0R039} & \texttt{1R133} & 243 & \texttt{0R164} & \texttt{1R150} & 327 & \texttt{0R215} & \texttt{1R220} \\
76 & \texttt{0R091} & \texttt{1R092} & 160 & \texttt{0R067} & \texttt{1R038} & 244 & \texttt{0R255} & \texttt{1R256} & 328 & \texttt{0R131} & \texttt{1R329} \\
77 & \texttt{0R021} & \texttt{1R087} & 161 & \texttt{0R175} & \texttt{0L024} & 245 & \texttt{0R257} & \texttt{1R244} & 329 & \texttt{0R217} & \texttt{1R330} \\
78 & \texttt{0R093} & \texttt{1R094} & 162 & \texttt{0R176} & \texttt{1R043} & 246 & \texttt{0R092} & \texttt{1R245} & 330 & \texttt{0R205} & \texttt{1R331} \\
79 & \texttt{0L095} & \texttt{1R061} & 163 & \texttt{0R177} & \texttt{1R178} & 247 & \texttt{0R169} & \texttt{1R258} & 331 & \texttt{0R217} & \texttt{1R332} \\
80 & \texttt{0L080} & \texttt{0L033} & 164 & \texttt{0R142} & \texttt{1R136} & 248 & \texttt{0R226} & \texttt{1R259} & 332 & \texttt{0R327} & \texttt{1R333} \\
81 & \texttt{1R096} & \texttt{1R081} & 165 & \texttt{0R179} & \texttt{1R150} & 249 & \texttt{0R205} & \texttt{1R260} & 333 & \texttt{0R092} & \texttt{1R296} \\
82 & \texttt{0L097} & \texttt{1L082} & 166 & \texttt{0R180} & \texttt{1R173} & 250 & \texttt{0R261} & \texttt{1R262} &  &  &  \\
83 & \texttt{1R098} & \texttt{0R083} & 167 & \texttt{0R181} & \texttt{1R182} & 251 & \texttt{0R263} & \texttt{1R140} &  &  &  \\
\bottomrule
\end{longtable}
\endgroup


\end{document}
